\documentclass[12pt]{article}

\setlength\parindent{0pt}
\newcommand{\myt}[1]{\textbf{\underline{#1}}}

\usepackage{mathtools}
\usepackage{amssymb}
\usepackage{tikz,ifthen,amsmath,amssymb,fancyhdr,comment,lastpage}

\title{\vspace{-15ex}Math 239 Lecture 15\vspace{-1ex}}
\date{June 10th, 2015}
\author{Graham Cooper}

\begin{document}
	\maketitle
	
	\section*{Counting Binary Trees}
	\begin{center}\begin{tikzpicture}[
		level distance=45 pt,
		every node/.style={circle,draw},
		level 1/.style={sibling distance=200 pt},
		level 2/.style={sibling distance=100 pt},
		level 3/.style={sibling distance=60 pt}
		]
		\node {}
		child {node {}
			child {node {}
				child {node {}}
				child {node {}}
			}
			child {node {}
				child [missing]
				child {node {}}
			}
		}
		child {node {}
			child [missing]
			child {node{}
				child {node {}}
				child {node {}}
			}	
		}
		;
	\end{tikzpicture}\end{center}

	\myt{Definition: } A binary tre is either:
	\begin{enumerate}
		\item an emtpy tree $\epsilon$
		\item a root vertex with a left branch anda right branch ech of which is a binary tree
	\end{enumerate}
	
	\subsection*{how many binary trees have n nodes?}
	for 3 nodes:\\
	\begin{center}\begin{tikzpicture}[
		level distance=45 pt,
		every node/.style={circle,draw},
		level 1/.style={sibling distance=200 pt},
		level 2/.style={sibling distance=100 pt},
		level 3/.style={sibling distance=60 pt}
		]
		\node {}
		child {node {}
		}
		child {node {}
		}
		;
	\end{tikzpicture}\end{center}
	\begin{center}\begin{tikzpicture}[
		level distance=45 pt,
		every node/.style={circle,draw},
		level 1/.style={sibling distance=200 pt},
		level 2/.style={sibling distance=100 pt},
		level 3/.style={sibling distance=60 pt}
		]
		\node {}
		child {node {}
			child {node {}}
		}
		child [missing]
		;
	\end{tikzpicture}\end{center}
\begin{center}\begin{tikzpicture}[
	level distance=45 pt,
	every node/.style={circle,draw},
	level 1/.style={sibling distance=200 pt},
	level 2/.style={sibling distance=100 pt},
	level 3/.style={sibling distance=60 pt}
	]
	\node {}
	child {node {}
		child [missing]
		child {node {}}
	}
	child [missing]
	;
	\end{tikzpicture}\end{center}
\begin{center}\begin{tikzpicture}[
	level distance=45 pt,
	every node/.style={circle,draw},
	level 1/.style={sibling distance=200 pt},
	level 2/.style={sibling distance=100 pt},
	level 3/.style={sibling distance=60 pt}
	]
	\node {}
	child[missing]
	child {node {}
		child {node {}}
		child[missing]
	}
	;
	\end{tikzpicture}\end{center}
	\begin{center}\begin{tikzpicture}[
	level distance=45 pt,
	every node/.style={circle,draw},
	level 1/.style={sibling distance=200 pt},
	level 2/.style={sibling distance=100 pt},
	level 3/.style={sibling distance=60 pt}
	]
	\node {}
	child[missing]
	child {node {}
		child[missing]
		child {node {}}
	}
	;
	\end{tikzpicture}\end{center}
	
	Let T be the set of all binary trees of any size. For each t $\in$ T, define weight functions w(t) to be the number of nodes in T. Let T(x) be the generating series fo T with respect to w. The answer to our question is $[x^n]t(x)$\\
	
	For any non-empty tree t, we can decompose it into 3 parts, root node r, the left tree $t_1$, the right tree $t_2$. $t_1$, $t_2 \in T$. This belongs to the cartesian product $r \times T \times T$\\
	
	Let S = $\{\epsilon \} \cup \{r \} \times T \times T$\\
	
	Then each tree t $\in$ T corresponds to one element of S in a bijection.\\
	
	\begin{center}\begin{tikzpicture}[
		level distance=45 pt,
		every node/.style={circle,draw},
		level 1/.style={sibling distance=200 pt},
		level 2/.style={sibling distance=100 pt},
		level 3/.style={sibling distance=60 pt}
		]
		\node {}
		child {node {}
			child {node {}
				child {node {}}
				child {node {}}
			}
			child {node {}}
		}
		child {node {}
			child {node {}}
			child {node {}}
		}
		;
	\end{tikzpicture}\end{center}
	
	becomes
	
	\begin{center}\begin{tikzpicture}[
		level distance=45 pt,
		every node/.style={circle,draw},
		level 1/.style={sibling distance=200 pt},
		level 2/.style={sibling distance=100 pt},
		level 3/.style={sibling distance=60 pt}
		]
		\node {0}
		;
	\end{tikzpicture}\end{center}
	\begin{center}\begin{tikzpicture}[
		level distance=45 pt,
		every node/.style={circle,draw},
		level 1/.style={sibling distance=200 pt},
		level 2/.style={sibling distance=100 pt},
		level 3/.style={sibling distance=60 pt}
		]
		\node {}
		child {node {}
			child {node {}
			}
			child {node {}}
		}
		child {node {}
		}
		;
		\end{tikzpicture}\end{center}
		\begin{center}\begin{tikzpicture}[
			level distance=45 pt,
			every node/.style={circle,draw},
			level 1/.style={sibling distance=200 pt},
			level 2/.style={sibling distance=100 pt},
			level 3/.style={sibling distance=60 pt}
			]
			\node {}
			child {node {}
			}
			child {node {}
			}
			;
		\end{tikzpicture}\end{center}
	
	$$T \rightleftarrows (r, t_1, t_2)$$
	$$w(r, t_1, t_2) = 1 + w(t_1) + w(t_2) = w(t)$$
	
	We can appy product lemma
	$$\Phi_S(x) = 1 + x \cdot t(x) \cdot t(x)$$
	Due to the bijection t(x) = $\Phi_S(x)$\\
	$$\implies t(x) = 1 + xt(x)^2$$	
	$$0 = xt(X)^2 - t(x) + 1$$
	$$= 4x(xt(x)^2 - t(x) + 1)$$
	$$= 4x^2t(x)^2 - 4xt(x) + 4x$$
	$$= (4x^2t(x)^2 - 4xt(x) + 1) - 1 + 4x$$
	$$= (2xt(x) - 1)^2 - 1 + 4x$$
	so:
	$$(2xt(x)-1)^2 = 1 - 4x = ((1-4x)^{1/2})^2$$
	from A2:
	$$2xt(x) - 1 = +/-(1-4x)^{1/2}$$
	$$= +/-(1-2\sum_{n\geq 0}\frac{1}{n+1}{2n \choose n}x^{n+1})$$
	Constant term is -1 so we pick - over +, which means:
	$$2xt(x) - 1 = -1 + 2 \sum_{n \geq 0}\frac{1}{n+1}{2n \choose n}x^{n+1}$$
	$$\implies t(x) = \sum_{n\geq 0}\frac{1}{n+1}{2n \choose n}x^n$$
	
	So the number of binary trees iwth n nodes is $\frac{1}{n+!}{2n \choose n}$\\
	(n = 3, $\frac{1}{4}{6 \choose 3} = \frac{1}{4}\frac{6 * 5 * 4}{3 * 2} = 5)$\\
	
	\myt{Catalan number: $\frac{1}{n+1}{2n \choose n}$}\\
	
	\section*{Graph Theory}
	
	\myt{Definition: } A graph G is a pair of sets (V, E) (or (V(G), E(G)) where V is a set of objects called vertices and E is a set of unordered pairs of V called edges.\\
	
	Example: Define G = (V,E) where\\
	V = \{a,b,c,d\}, E = \{\{a,b\}, \{b,c\}, \{c,d\}, \{a,d\} \}\\
	
	Graphical Representation of G.\\
	See notes.\\
	
	In graph theory, we mainly care about the "structure" of the graphs, e.g. what are the vertices, which pairs for edges\\
	
	Terminologies:
	\begin{itemize}
		\item Two vertices u,v are adjacent if \{u,v\} is an edge Example above, a is adjacent to b and d but not c
		\item If u is adjacent to v, then u is a neighbour of v, the set of all neighbours of u is the neighbourhood, denoted by N(v). Example: N(a) = \{b,d\}
		\item An edge, e = \{u,v\} is incident with u and v, e joins u and v
	\end{itemize}
	
	
	
\end{document}
